\section{ScrapeBench: Experimental Setup}
\label{sec:experiments}

\subsection{Implementation and independent interventions}
ScrapeBench is a resumable Python harness built on \texttt{diffusers},
\texttt{transformers}, \texttt{insightface}, and Pillow. Every condition starts
from the same T0 gallery; T1--T6 are independent interventions rather than a
seven-stage cascade. We evaluate T0 (unaltered), T1 (JPEG Q70, 4:2:0), T2
(360-CSS-pixel browser raster at DPR 1), T3 (face crop followed by WebP 80), T4
(Real-ESRGAN x2 and JPEG 92), T5 (SD1.5 img2img at strength .30), and T6 (50-step
diffusion purification). Clean and protected branches share the loader, crop,
intervention, attacker, prompt, seed, and scorer.

\subsection{Data, defenses, and attackers}
The fixed CelebA-HQ-100 protocol contains 100 identity-labelled subjects, each
with eight gallery images and two held-out references. The public split is 10
development, 80 core-test, and 10 untouched audit identities; hyperparameters
are frozen on development before the 90 non-development identities are run.
The primary comparison uses the clean control, official GAP-Diff and
IDProtector checkpoints, and ARCP-v2. Each published baseline is run at its
official released setting and its measured pixel budget is reported, rather
than silently forcing methods to a nominal but implementation-dependent
budget. ARCP-v2 applies its bound in 8-bit pixel space, with
$\epsilon_{\min}=2/255$, $\epsilon_{\max}=16/255$, a $64^2$ band-limited grid,
200 optimization steps, reader weight $.75$, and margin $.35$. The older
40-identity FSMG/ASPL-lite experiment is retained only as a legacy diagnostic
and is not mixed with the revised tables.

Stage A tests two zero-shot readers: IP-Adapter-FaceID on SD1.5 and InstantID on
SDXL, both with 30 denoising steps and guidance 5.0. Stage B trains the SD1.5
UNet for 1,000 steps on each acquired gallery and generates 30 images for each
of four preregistered Web scenarios: social-media avatar, smartphone selfie,
low-light indoor social photograph, and outdoor travel photograph. A separate
VGGFace2/JPEG-Q70 block preserves GAP-Diff's four original prompts for direct
comparison with its Table I and Figure 8 protocol.

\subsection{Metrics and statistical contract}
The four headline output metrics are face-detection failure rate (FDFR),
identity similarity (ISM), BRISQUE, and SER-FIQ. We additionally report the
same-identity clean-minus-protected ISM gap, its identity-bootstrap 95\% CI,
retention $\mathrm{RR}(T)=\Delta(T)/\Delta(T0)$ with a bootstrap CI, change from
T0, Holm-adjusted paired tests, and a TOST equivalence test at the preregistered
0.03 ISM bound. Non-significance is not described as equivalence or collapse.
The primary tables will be generated only from the append-only revised run log;
legacy and revised rows are never pooled.

\subsection{Execution status}
At manuscript revision time the 100-identity dataset and split are frozen, and
the official GAP-Diff and ARCP-v2 protected galleries are complete for all 90
evaluation identities. IDProtector and the downstream Stage-A and Stage-B
evaluations are still running. Consequently, this draft intentionally
contains no revised performance numbers or superiority claim; Section~\ref{sec:results}
defines the tables that will be populated from the completed run.
